\documentclass[../../main.tex]{subfiles}
\begin{document}

\paragraph{【新課程】}
{\bf [解]}

\begin{enumerate}[(1)]
  \item 直線$l$の方向ベクトルを$\vec{l}$，直線$m$の方向ベクトルを$\vec{m}$とすると，
    \begin{align}
      \vec{l} &=
      \begin{pmatrix} 2-1 \\ 1-1 \\ 1-0
      \end{pmatrix} \nonumber\\
      &=
      \begin{pmatrix} 1 \\ 0 \\ 1
      \end{pmatrix} \\
      \vec{m} &=
      \begin{pmatrix} 1-1 \\ 3-1 \\ 2-1
      \end{pmatrix} \nonumber\\
      &=
      \begin{pmatrix} 0 \\ 2 \\ 1
      \end{pmatrix}
    \end{align}
    である．求める平面$\pi$は直線$l$を含むので点$(1,1,0)$を通り，かつ$\vec{l}$および$\vec{m}$に平行である．
    平面$\pi$の法線ベクトルを $\vec{n}_0 =
    \begin{pmatrix} a \\ b \\ c
    \end{pmatrix} \neq \vec{0}$ とおくと，
    \begin{align}
      \vec{n}_0 \cdot \vec{l} = a + c = 0 \implies c = -a, \qquad
      \vec{n}_0 \cdot \vec{m} = 2b + c = 0 \implies 2b = -c = a
    \end{align}
    であるから，$a = -2$ とおくと $b = -1$，$c = 2$ となり，法線ベクトルとして $\vec{n}_0 =
    \begin{pmatrix} -2 \\ -1 \\ 2
    \end{pmatrix}$ がとれる．
    したがって，点$(1,1,0)$を通る平面$\pi$の方程式は
    \begin{align}
      -2(x-1) - (y-1) + 2(z-0) = 0 \iff -2x - y + 2z + 3 = 0 \quad \cdots\text{（答）}
    \end{align}
    である（$2x + y - 2z - 3 = 0$ と表してもよい）．

  \item 直線$l$上の点$P$は実数$s$を用いて
    \begin{align}
      P(1+s, 1, s)
    \end{align}
    と表せ，直線$m$上の点$Q$は実数$t$を用いて
    \begin{align}
      Q(1, 1+2t, 1+t)
    \end{align}
    と表せる．点$R(2,0,1)$を通り$l$，$m$と交わる直線を$n$とすると，交点はそれぞれ$P$，$Q$であり，3点$R, P, Q$は一直線上にある．
    すなわち，$\overrightarrow{RP}$ と $\overrightarrow{RQ}$ は平行である．ここで
    \begin{align}
      \overrightarrow{RP} &=
      \begin{pmatrix} 1+s-2 \\ 1-0 \\ s-1
      \end{pmatrix} \nonumber\\
      &=
      \begin{pmatrix} s-1 \\ 1 \\ s-1
      \end{pmatrix} \\
      \overrightarrow{RQ} &=
      \begin{pmatrix} 1-2 \\ 1+2t-0 \\ 1+t-1
      \end{pmatrix} \nonumber\\
      &=
      \begin{pmatrix} -1 \\ 1+2t \\ t
      \end{pmatrix}
    \end{align}
    である．$\overrightarrow{RP}$の$x$成分と$z$成分が等しいことから，$\overrightarrow{RQ}$の$x$成分と$z$成分も等しくなければならない．
    ゆえに
    \begin{align}
      -1 = t \implies t = -1
    \end{align}
    このとき $\overrightarrow{RQ} =
    \begin{pmatrix} -1 \\ -1 \\ -1
    \end{pmatrix}$ となる．
    $\overrightarrow{RP} = k \overrightarrow{RQ}$ とおくと，第2成分を比較して
    \begin{align}
      1 = -k \implies k = -1
    \end{align}
    であるから，$\overrightarrow{RP} =
    \begin{pmatrix} 1 \\ 1 \\ 1
    \end{pmatrix}$ となり，
    \begin{align}
      s - 1 = 1 \implies s = 2
    \end{align}
    を得る．
    よって，求める交点の座標は
    \begin{align}
      l \text{ と } n \text{ の交点: } (3, 1, 2), \qquad
      m \text{ と } n \text{ の交点: } (1, -1, 0) \quad \cdots\text{（答）}
    \end{align}
    である．
\end{enumerate}

\paragraph{【旧課程】}
{\bf [解]}

\begin{enumerate}[(1)]
  \item $\cos 2x = 2\cos^2 x - 1$ より，
    \begin{align}
      f(x) &= -\cos x - \dfrac{a}{4}(2\cos^2 x - 1) \nonumber\\
      &= -\dfrac{a}{2}\cos^2 x - \cos x + \dfrac{a}{4}
    \end{align}
    である．$t = \cos x$ とおくと，$0 \leqq x \leqq 2\pi$ より $-1 \leqq t \leqq 1$ である．
    $t$ の関数 $g(t)$ を
    \begin{align}
      g(t) &= -\dfrac{a}{2}t^2 - t + \dfrac{a}{4} \nonumber\\
      &= -\dfrac{a}{2}\left( t + \dfrac{1}{a} \right)^2 + \dfrac{1}{2a} + \dfrac{a}{4}
    \end{align}
    とおく．放物線 $y = g(t)$ は上に凸で，軸は直線 $t = -\dfrac{1}{a} < 0$ である．
    \begin{enumerate}[(i)]
      \item $0 < a \leqq 1$ のとき：
        $-\dfrac{1}{a} \leqq -1$ であるから，$g(t)$ は $-1 \leqq t \leqq 1$ で単調に減少する．
        したがって，
        \begin{itemize}
          \item 最大値は $t = -1$（$x = \pi$）のとき，
            \begin{align}
              g(-1) = -\dfrac{a}{2} + 1 + \dfrac{a}{4} = 1 - \dfrac{a}{4}
            \end{align}
          \item 最小値は $t = 1$（$x = 0, 2\pi$）のとき，
            \begin{align}
              g(1) = -\dfrac{a}{2} - 1 + \dfrac{a}{4} = -1 - \dfrac{a}{4}
            \end{align}
        \end{itemize}
      \item $a > 1$ のとき：
        $-1 < -\dfrac{1}{a} < 0$ であるから，頂点 $t = -\dfrac{1}{a}$ は $-1 \leqq t \leqq 1$ の範囲に含まれる．
        また，区間の両端における値の差は
        \begin{align}
          g(-1) - g(1) = \left( 1 - \dfrac{a}{4} \right) - \left( -1 - \dfrac{a}{4} \right) = 2 > 0
        \end{align}
        より，常に $g(-1) > g(1)$ である．
        したがって，
        \begin{itemize}
          \item 最大値は $t = -\dfrac{1}{a}$（$\cos x = -\dfrac{1}{a}$ を満たす $x$）のとき，
            \begin{align}
              g\left( -\dfrac{1}{a} \right) = \dfrac{1}{2a} + \dfrac{a}{4}
            \end{align}
          \item 最小値は $t = 1$（$x = 0, 2\pi$）のとき，
            \begin{align}
              g(1) = -1 - \dfrac{a}{4}
            \end{align}
        \end{itemize}
    \end{enumerate}
    以上をまとめると，
    \begin{align}
      \begin{cases}
        0 < a \leqq 1 \text{ のとき} & \text{最大値 } 1 - \dfrac{a}{4}, \quad \text{最小値 } -1 - \dfrac{a}{4} \\
        a > 1 \text{ のとき} & \text{最大値 } \dfrac{1}{2a} + \dfrac{a}{4}, \quad \text{最小値 } -1 - \dfrac{a}{4}
      \end{cases}
      \quad \cdots\text{（答）}
    \end{align}
    である．

  \item $f(x)$ を$x$で微分すると，
    \begin{align}
      f'(x) &= \sin x + \dfrac{a}{2}\sin 2x = \sin x (1 + a\cos x) \\
      f''(x) &= \cos x + a\cos 2x = 2a\cos^2 x + \cos x - a
    \end{align}
    である．$t = \cos x$（$-1 \leqq t \leqq 1$）とおき，
    \begin{align}
      h(t) = 2at^2 + t - a
    \end{align}
    とする．$2$次方程式 $h(t) = 0$ の判別式は $D = 1 + 8a^2 > 0$ であるから，相異なる$2$つの実数解
    \begin{align}
      \alpha = \dfrac{-1 - \sqrt{1+8a^2}}{4a}, \qquad \beta = \dfrac{-1 + \sqrt{1+8a^2}}{4a} \quad (\alpha < \beta)
    \end{align}
    をもつ．
    ここで
    \begin{align}
      h(0) = -a < 0, \qquad h(1) = a + 1 > 0
    \end{align}
    であるから，常に $0 < \beta < 1$ である．
    一方，$\alpha$ については
    \begin{align}
      h(-1) &= 2a(-1)^2 + (-1) - a \nonumber\\
      &= a - 1
    \end{align}
    であるから，$a$ の値によって次のように場合分けされる．
    \begin{enumerate}[(i)]
      \item $0 < a < 1$ のとき：
        $h(-1) = a - 1 < 0$ より $\alpha < -1$ である．
        したがって，$-1 \leqq t \leqq 1$ において $h(t) = 0$ となるのは $t = \beta$ のみである．
        $0 < \beta < 1$ より，$\cos x = \beta$ を満たす $x \in (0, 2\pi)$ は
        \begin{align}
          x_1 \in \left( 0, \dfrac{\pi}{2} \right), \qquad x_4 = 2\pi - x_1 \in \left( \dfrac{3\pi}{2}, 2\pi \right)
        \end{align}
        の$2$個存在する．
        $t = \cos x$ は $0 \leqq x \leqq \pi$ で単調減少，$ \pi \leqq x \leqq 2\pi$ で単調増加であるから，
        \begin{itemize}
          \item $0 \leqq x < x_1$ および $x_4 < x \leqq 2\pi$ では $\cos x > \beta$ より $h(t) > 0$，すなわち $f''(x) > 0$（下に凸）
          \item $x_1 < x < x_4$ では $\cos x < \beta$ より $h(t) < 0$，すなわち $f''(x) < 0$（上に凸）
        \end{itemize}
        となり，$x = x_1, x_4$ の前後で $f''(x)$ の符号が変化する．したがって，変曲点は$2$個である．

      \item $a = 1$ のとき：
        $\alpha = -1$，$\beta = \dfrac{1}{2}$ である．
        $h(t) = 2(t+1)\left( t - \dfrac{1}{2} \right)$ と因数分解される．
        $\cos x = \dfrac{1}{2}$ となる点は $x_1 = \dfrac{\pi}{3}$，$x_4 = \dfrac{5\pi}{3}$ であり，
        $\cos x = -1$ となる点は $x = \pi$ である．
        $x \neq \pi$ では常に $t+1 = \cos x + 1 > 0$ であるから，$x = \pi$ の前後で $f''(x)$ の符号は変化しない（$f''(\pi)=0$ ではあるが変曲点ではない）．
        $x_1, x_4$ の前後でのみ $f''(x)$ の符号が変化するから，変曲点は$2$個である．
        凹凸の変化は (i) と同様である．

      \item $a > 1$ のとき：
        $h(-1) = a - 1 > 0$ より $-1 < \alpha < 0$ である．
        したがって，$-1 < \alpha < 0 < \beta < 1$ となり，$-1 \leqq t \leqq 1$ において $h(t) = 0$ となる解が $t = \alpha, \beta$ の$2$つ存在する．
        これらに対応する $x \in (0, 2\pi)$ を小さい順に
        \begin{align}
          x_1 &= \arccos\beta, &x_2 &= \arccos\alpha, \nonumber\\
          x_3 &= 2\pi - x_2, &x_4 &= 2\pi - x_1
        \end{align}
        とおくと，$0 < x_1 < \dfrac{\pi}{2} < x_2 < \pi < x_3 < \dfrac{3\pi}{2} < x_4 < 2\pi$ である．
        各区間における $f''(x)$ の符号は，
        \begin{itemize}
          \item $0 \leqq x < x_1$：$\cos x > \beta$ より $h(t) > 0$，すなわち $f''(x) > 0$（下に凸）
          \item $x_1 < x < x_2$：$\alpha < \cos x < \beta$ より $h(t) < 0$，すなわち $f''(x) < 0$（上に凸）
          \item $x_2 < x < x_3$：$\cos x < \alpha$ より $h(t) > 0$，すなわち $f''(x) > 0$（下に凸）
          \item $x_3 < x < x_4$：$\alpha < \cos x < \beta$ より $h(t) < 0$，すなわち $f''(x) < 0$（上に凸）
          \item $x_4 < x \leqq 2\pi$：$\cos x > \beta$ より $h(t) > 0$，すなわち $f''(x) > 0$（下に凸）
        \end{itemize}
        となり，4点 $x_1, x_2, x_3, x_4$ のすべてにおいて $f''(x)$ の符号が変化する．したがって，変曲点は$4$個である．
    \end{enumerate}

    まとめると，
    \begin{itemize}
      \item $0 < a \leqq 1$ のとき：
        変曲点の個数は$2$個．
        凹凸の変化は，
        $0 \leqq x < x_1$ で下に凸，
        $x_1 < x < x_4$ で上に凸，
        $x_4 < x \leqq 2\pi$ で下に凸
        （ただし $x_1, x_4$ は $\cos x = \dfrac{\sqrt{1+8a^2}-1}{4a}$ を満たす $0 < x_1 < \dfrac{\pi}{2} < x_4 < 2\pi$）．
      \item $a > 1$ のとき：
        変曲点の個数は$4$個．
        凹凸の変化は，
        $0 \leqq x < x_1$ で下に凸，
        $x_1 < x < x_2$ で上に凸，
        $x_2 < x < x_3$ で下に凸，
        $x_3 < x < x_4$ で上に凸，
        $x_4 < x \leqq 2\pi$ で下に凸
        （ただし $x_1, x_4$ は $\cos x = \dfrac{\sqrt{1+8a^2}-1}{4a}$，
        $x_2, x_3$ は $\cos x = -\dfrac{\sqrt{1+8a^2}+1}{4a}$ を満たす $0 < x_1 < x_2 < \pi < x_3 < x_4 < 2\pi$）．
    \end{itemize}
    $\cdots\text{（答）}$
\end{enumerate}

\end{document}
